% Folder 34, batch 03 — archivists' pages 41–60.
% Pass: Opus 5.5 (claude-opus-5-5), 2026-10-03 — first pass, unchecked against the pages by a human.
%
% Page 51 is blank and carries no \page{}; its absence is the record.
\documentclass[11pt,a4paper]{article}
\input{../preamble/grothendieck}

\folder{34}
\batch{3}
\pages{41}{60}
\foldertitle{SGA 7 : notes manuscrites (s.d.).}
\dating{[vers 1967-1973]}
\watermark{Édition de démonstration}

\begin{document}

\page{41}
\note{p.~40 de l'auteur. La page s'ouvre au milieu d'une démonstration commencée avant ce lot.}

\note{Le haut de la page est un passage encadré, barré de deux longs traits obliques.}
\struck{\ill{} isom. canonique $\pi_1(S,\xi) \simeq \pi_1(\operatorname{Spec} k, \bar k) = \operatorname{Gal}(\bar k/k)$. De plus, l'hom. canonique $\pi_1(U,\xi) \to \pi_1(S,\xi) = \pi_0$ (déf.) est surjectif, et \ill{} \ill{} à la structure suivante :}

\uncertain{Quotientant} $E$ par son sous-groupe de torsion (ce qui ne change pas la conclusion, ni l'hyp.), on se ramène \uncertain{par ex.} au cas où $E$ est \ill{} sur $\mathbf{Z}_\ell$, l'hyp. d'action \uncertain{unipotente} signifiant maintenant que l'image de $I$ dans $\operatorname{Aut}(E_0)$ ($E_0 = E \otimes_{\mathbf{Z}_\ell} \mathbf{Z}/\ell\mathbf{Z}$) est d'ordre premier à $p$. Mais \ill{} \ill{} \ill{} $I$ est \ill{} \uncertain{limite projective} des \ill{} $I_\alpha$, \ill{} l'image de $I$ dans $\operatorname{Aut}(E_0)$ est \ill{} l'intersection des \ill{} \ill{} $I_\alpha$ \ill{} \ill{} \ill{} premier à $p$. Donc $I' = I_t$ \ill{} \ill{} d'un groupe \ill{} \ill{} \ill{} premier à $p$ \ill{}

\uncertain{(B)} On utilise maintenant le \ill{} de structure \uncertain{suivant}, qui est une des variantes du lemme d'Abhyankar :

\marginal{devrait \ill{} avant l'actuel n°~2}

\textbf{Lemme 8.5.} Soient $A$ un anneau local noeth. régulier, $k$ son corps résiduel, $p$ l'exp. car. de $k$, $(x_i)_{1 \leq i \leq n}$ une partie d'un syst. régulier

\page{42}
\note{p.~41 de l'auteur.}
de paramètres de $A$, $x = \prod x_i$, \struck{\ill{}} $S = \operatorname{Spec} A$, $U = \operatorname{Spec} A_x = S_x$, $\xi$ un pt géom. de $U$.

a) $\xi$ définit une clôture \uncertain{séparable} $\bar k$ de $k$ et on a un isom. canonique
\[
(8.5.1) \qquad \pi_1(S,\xi) \simeq \pi_1(\operatorname{Spec} k, \bar k) = \operatorname{Gal}(\bar k/k).
\]

b) L'hom. canonique
\[
(8.5.2) \qquad \pi_1(U,\xi) \longrightarrow \pi_1(S,\xi) \simeq \operatorname{Gal}(\bar k/k)
\]
est surjectif, et définit donc $\pi_1(U,\xi)$ comme extension de $\operatorname{Gal}(\bar k/k)$ par un groupe $I$ ; \struck{\ill{} \ill{} \ill{} : $\pi_1(\tilde S, \tilde\xi)$, où $\tilde S$ \ill{} le hensélisé strict}
\[
(8.5.3) \qquad 1 \longrightarrow I \longrightarrow \pi_1(U,\xi) \longrightarrow \operatorname{Gal}(\bar k/k) \longrightarrow 1 ,
\]
et on a un isom. can.
\[
(8.5.4) \qquad I \simeq \pi_1(\tilde U, \tilde\xi) ,
\]
où $\tilde S$ désigne le hensélisé strict de $S$ relativement à $\bar k/k$, \struck{\ill{}} $\tilde U$ l'image inverse de $U$ dans $\tilde S$, et $\tilde\xi$ un \ill{} de $\xi$ à $\tilde S$.

c) Soit $I_t$ le plus grand quotient de $I$ qui soit \ill{} \ill{} premier à $p$. Alors on a un isom. canonique
\[
(8.5.5) \qquad I_t \simeq \Bigl(\prod_{\ell \neq p} \mathbf{Z}_\ell(1)\Bigr)^n ,
\]
\[
\mathbf{Z}_\ell(1) \simeq T_\ell(k(\xi)^*) \simeq T_\ell(\bar k^*) .
\]
\note{Dans la première formule, un premier membre $\bigl(\hat{\mathbf{Z}}'(1)\bigr)^n$ est biffé et remplacé par le produit ; la seconde portait le numéro (8.5.6), biffé. Une ligne biffée suit : $\hat{\mathbf{Z}}'(1) = \prod \mathbf{Z}_\ell(1)$.}

\marginal{ne sert pas dans la \uncertain{démonstration} de 8.1.}
Pour un revêtement principal \struck{\ill{}} \ill{} $S'$ de $S$, de groupe $\Gamma$, la représentation correspondante $\pi \to \Gamma$, i.e. un hom. $I \to \Gamma$, \struck{qui} se factorise par $I_t$ si et seulement si \struck{\ill{}}

\page{43}
\note{p.~42 de l'auteur.}
\ill{} la suite, s'il est modérément ramifié le long des \struck{hypersurfaces \ill{} des $x_i$} \add{diviseurs $\operatorname{div}(x_i)$}.

d) L'action naturelle de $\pi_1(U,\xi)$ sur \struck{$I_t =$} $\prod_{\ell \neq p} I/[I,I]^{(\ell)}$, qui se factorise \struck{\ill{}} évidemment par une action de $\operatorname{Gal}(\bar k/k)$, s'identifie via l'isomorphisme (8.5.5) \struck{\ill{} \ill{}} au \uncertain{produit} de l'action naturelle de $\operatorname{Gal}(\bar k/k)$ sur $T_\ell(\bar k^*)$, \struck{\ill{}} i.e. les opérations de $\operatorname{Gal}(\bar k/k)$ sur les ${}_{\ell^\infty}(\bar k^*) \subset \bar k^*$.

Démonstration laissée au lecteur ---

\textbf{Corollaire 8.6.} Sous les conditions de 8.5, supposons qu'aucune extension finie de $k$ ne contienne \ill{} les racines \ill{} de $1$, pour tout entier $\nu$. Alors \ill{} \ill{} représentation $\ell$-adique de $\pi_1(U,\xi)$ \add{\uncertain{modérée}, i.e. telle que} l'action induite de $I$ \add{se factorise par une action} de $I_t$, \struck{soit quasi-unipotente.} cette dernière est \emph{quasi-unipotente}.

\textbf{Remarque 8.7.} \struck{$t = u \prod x_i^{\alpha_i}$ \ill{} \ill{} $\alpha_i \geq 1$, et $x_i \in \mathfrak{m}_V \subset \mathfrak{m} = tV$} \struck{$u_i \in A^* \subset V^*$} Dans le cas d'application de 8.6, \ill{} aura \struck{\ill{}} $n = 1$ (un seul $x_i$, soit $x_1$) et $\alpha_1 = 1$. D\supplied{onc} \struck{\ill{} le cas \ill{}, la structure du groupe} $n = 1$, et l'hom. $I \to I'$ induit un isom.
\[
I_t \xrightarrow{\ \sim\ } I'_t .
\]

\textbf{Remarque 8.8.} La démonstration de 8.1 \ill{} \ill{}, \uncertain{hyp.} \ill{} des \uncertain{propriétés} sur $X_K$, si \ill{} \ill{} sait que l'on peut \ill{} un \ill{} \ill{} $V_\alpha$ des pts \uncertain{géométriques} $y_\alpha$ de $S$ \ill{}

\page{44}
\note{p.~43 de l'auteur.}
au dessus duquel \add{$E_\alpha =$} $R^i f_{\alpha *}(\mathbf{Z}_{\ell\, X_\alpha})$ est constant tordu. En effet, le fait que \struck{$E$} $R^i f_*(\mathbf{Z}_{\ell\, X_K})$, \ill{} \ill{} \ill{} \ill{} \struck{$Y_\alpha$} \struck{$\to$} $Y_\alpha \to S_\alpha$ \uncertain{signifie} que $R^i f_{\alpha *}(\mathbf{Z}_{\ell\, X_\alpha})$ est \ill{} \ill{} \ill{} \ill{} de base $Y_\alpha \to S_\alpha$ et $Y \to Y_\alpha$. \ill{} \ill{} est trivial, par le \ill{} \ill{} de base \ill{} \ill{} \ill{} \ill{} \ill{} \ill{} \ill{} (SGA~4~XVI~1.6). \struck{D'autre} \struck{\ill{}, l'exis\ill{}} D'autre part, l'existence d'un $V_\alpha$ comme dessus est \ill{} \ill{} \ill{} \ill{} de la résolution des singularités \add{(sous la forme de Hironaka)} pour les schémas de type \uncertain{fini} et de dim. \ill{} $\leq \dim X_K \; \ill{}$ \uncertain{sur} $K_\alpha$ (SGA~5~II~5.1). \add{Cette hypothèse est \uncertain{vérifiée}} \struck{\ill{}} en particulier, par Hironaka, si $K$ est de car. nulle, et \struck{\ill{} \ill{} \ill{} \ill{} \ill{} \ill{} lorsque \ill{} par Abhyankar, si $\dim X_K \leq 2$.}

\note{Ce qui suit, jusqu'au corollaire 8.9 inclus, est encadré et barré en croix.}
\struck{(\ill{} $X$ de type \uncertain{fini} sur $X_K$.) \textbf{Corollaire 8.9.} Soit $V$ comme dans 8.1, \struck{\ill{} $K = \ill{}$} \ill{} \ill{} \ill{} \uncertain{nulle}, ou $\dim X_K \leq 2$. \struck{Soit $X_K$ \ill{} de \ill{} sur $K$,} pour $i \in \mathbf{Z}$, \struck{\ill{}} \ill{} \ill{} pour $V_\alpha$ du \uncertain{résultat} \ill{} \ill{} \ill{} \ill{} (\ill{} dans 8.2 a)) ou l'action de $I$ dans $H^i(X_{\bar K}, \mathbf{Z}_\ell)$ modérée (cas de 8.2 b) par exemple). Alors l'\add{action} \struck{\ill{}} de $I$ sur $H^i(X_{\bar K}, \mathbf{Z}_\ell)$ est quasi-unipotente.}
\marginal{\struck{\ill{} $K$ de \ill{}}}

\textbf{Corollaire 8.9.} Les conclusions de 8.2 restent \uncertain{vraies} dans le \uncertain{cas} où $X$ est \uncertain{seulement} supposé de t.f. sur $K$ \uncertain{(cas \ill{})}, pourvu que dans le cas a), on suppose \add{la caractéristique de $K$ égale à $0$,} \emph{ou} $\dim X_K \leq 2$.

\page{45}
\note{p.~44 de l'auteur.}
\section*{9. Réductions au cas projectif et lisse.}

9.0. Nous \ill{} \ill{}, la question de savoir, pour \add{$V$, $\alpha$} une situation donnée, si le th. de monodromie 7.1. est vrai pour tout $X_K$ de t.f. sur $K$, \struck{de dim} tel que
\[
\dim X_K \leq n .
\]
Nous allons d'ailleurs envisager la conjecture correspondante pour les $\pi$-modules
\[
H^i_!(X_{\bar K}, \mathbf{Q}_\ell) \qquad (\text{pour } X_K \text{ \emph{régulier}}).
\]
\add{\uncertain{la conjecture est vérifiée} pour les dim. $\leq n-1$.}
Nous supposerons que la conj. \ill{} \ill{}

9.1. \emph{Cas $X$ lisse.} \add{\emph{Relation entre $H^i$ et $H^i_!$.}} Dans ce cas, \struck{\ill{}} \add{et si $X$ est \uncertain{propre},} on a un isom. canonique
\[
H^i(X_{\bar K}, \mathbf{Q}_\ell) \simeq H^{2n-i}_!(X_{\bar K}, \mathbf{Q}_\ell)'(-n) ,
\]
\note{Dans la formule, un facteur $\otimes \mathbf{Q}_\ell(n)$ est raturé et remplacé par « $(-n)$ » ; un « $(=n)$ » reste lisible dans la rature.}
où $\mathbf{Q}_\ell(n)$ désigne le \uncertain{twist} de Tate $\otimes_{\mathbf{Q}_\ell} \mathbf{Q}_\ell(n)$, et $'$ désigne le dual. Cet isom. est compatible avec l'action de $\pi$ et \emph{a fortiori} de $I$. Comme $I$ opère trivialement sur $\mathbf{Q}_\ell(-n)$, il s'ensuit que l'action de $I$ sur $H^i(X_{\bar K}; \mathbf{Q}_\ell)$ est \uncertain{quasi-unip.} (resp. unip.) ssi \ill{} l'action sur $H^i_!(X_{\bar K}, \mathbf{Q}_\ell)$ l'est.\note{L'indice inférieur de ce $H$ est surchargé ; « ! » est une lecture.}

9.2. \add{\emph{Cas} \ill{} $H^i_!(X_{\bar K}, \mathbf{Q}_\ell)$, $X_K$ \uncertain{sép.} de t.f. sur $K$.} \struck{La suite \ill{} \ill{}} \struck{Pour} Soient $U$ un ouvert \struck{\ill{}} \add{dense} de $X_K$, $Y$ \ill{} complémentaire. Alors \ill{} \ill{} la suite exacte
\[
\cdots \to H^i_!(Y_{\bar K}, \mathbf{Q}_\ell) \to H^i_!(X_{\bar K}, \mathbf{Q}_\ell) \to H^i_!(U_{\bar K}, \mathbf{Q}_\ell) \to H^{i+1}_!(U_{\bar K}, \ldots
\]
\note{L'argument du premier terme est raturé et peu lisible ; $Y$ est une lecture. L'ordre des termes est celui de la page.}
qui montre, moyennant l'hyp. de récurrence, que

\page{46}
\note{p.~45 de l'auteur.}
l'action de $I$ sur $H^i_!(X_{\bar K})$ est quasi-unipotente ssi son action sur $H^i_!(U_{\bar K})$ l'est : \ill{} \ill{} \uncertain{birationnel} de la question.

Si donc on peut trouver un \struck{\ill{}} \add{\ill{} propre et lisse $\hat U_K$} de $U_K$, \ill{} \ill{} ramené au cas d'un schéma \emph{lisse} et \emph{propre}. Si on admet la résolution \add{(\uncertain{pour} \ill{})} des singularités \struck{\ill{}} \add{pour la clôture \uncertain{parfaite} de $K$,} pour les schémas de t.f. de dim. $\leq n$, on \ill{} \add{(\uncertain{pourvu} \ill{} $U$ affine)} \ill{}, \uncertain{sous} la condition que l'hypothèse de récurrence \ill{} \ill{} \ill{} toute extension \add{radicielle} finie \struck{\ill{}} $K'$ de $K$ (\ill{} \uncertain{remplaçant} la clôture \ill{} $V'$ de $V$ dans $K'$).

\textbf{Remarque 9.3.} \struck{D\ill{}} \ill{} \ill{} \ill{} de 7.1, ça \ill{} aussi pour $H^i(X_{\bar K}, \mathbf{Q}_\ell)$, si $X$ lisse \emph{séparé}. Le cas non \uncertain{séparé} s'y ramène, par la suite spectrale de Leray d'un recouvrement fini de $X$ par des ouverts affines.

\note{Le bas de la page est encadré et barré de traits obliques.}
\struck{9.3. Cas de $H^i_!(X_{\bar K}, \mathbf{Q}_\ell)$ pour $X$ \add{séparé de t.f.} quelconque. \struck{Soit \ill{}, On peut \ill{}} Si $U$ est un ouvert \ill{}, \ill{} $X$ admet \ill{}, et \ill{} \ill{} \ill{} un \ill{} \ill{} \ill{} de $X$ par $U$. On peut trouver \ill{} $U$ tel que \ill{} $U$ \ill{} \ill{} \ill{} soit \ill{}. \ill{} \ill{} \ill{} \ill{} une extension \add{finie} \ill{} radicielle $K'$ de $K$, telle que \ill{} $(U_{K'})_{\mathrm{red}}$ soit \emph{lisse} sur $K'$. On \ill{} \ill{} \ill{} \ill{} \ill{} de 9.2.}

9.3. \emph{Cas de} $H^i(X_{\bar K}, \mathbf{Q}_\ell)$ \emph{pour $X$ \uncertain{de} t.f. quelconque.}
\note{Ce dernier intitulé est tracé sur le cadre du passage barré ; il n'est pas certain qu'il soit conservé.}

\page{47}
\note{p.~46 de l'auteur.}
Cela \uncertain{prouve} :

\textbf{Proposition 9.4.} \struck{Pour} a) \supplied{Pour} la validité du th. de monodromie \add{(pour les $H^i_!(X_{\bar K}, \mathbf{Q}_\ell)$)}, pour $V$ fixé, et des $X_K$ avec $\dim X_K \leq n$, et tout $i$, il suffit de savoir qu'il est vrai pour tout $X_{K'}$, où $X$ est lisse \add{de dim.\ $\leq n$} sur \struck{$K'$} \add{une} extension finie radicielle $K'$ de $K$, considérée comme corps des fractions de la clôture normale $V'$ de $V$ dans $K'$. \struck{\ill{}} On peut \struck{\ill{} \ill{}} \ill{} \ill{} supposer les $X_{K'}$ affines. \struck{a)} b) Dans le cas où \struck{la résolution des singularités \ill{}} toute extension de t.f. et de degré de transcendance $\leq n$ sur la clôture parfaite $K^{p^{-\infty}}$ de $K$ admet un modèle projectif et lisse, il suffit de supposer \add{les} $X_{K'}$ projectifs. c) Enfin, \struck{si} ces conditions \struck{préalables} \ill{} étant vérifiées, \ill{} a) \ill{} vérifiée, donc pour tout $X_K$ de t.f. et \emph{lisse} sur $K$, \add{et de dim.\ $\leq n$,} le th. de monodromie est vrai aussi pour les $H^i_*(X_{\bar K}, \mathbf{Q}_\ell)$.

9.5. Nous allons maintenant examiner le th. de monodromie pour $H^i(X_{\bar K}, \mathbf{Q}_\ell)$, \add{pour} $X_K$ de t.f., en supposant que \struck{les conditions de 7.1 \ill{} \ill{} vérifiées \ill{} tout \ill{}.} \struck{On procède par récurrence sur $i$, le cas $i < 0$ étant trivial.} \ill{} \uncertain{prouvé} pour les $H^j(X_{\bar K}, \mathbf{Q}_\ell)$ pour $j < i$.

\page{48}
\note{p.~47 de l'auteur.}
On peut supposer \ill{} \ill{} $U$ affine, donc quasi-projectif. \struck{\ill{}} On admet que la résolution des singularités sur la clôture parfaite \struck{\ill{}} de $K$ \ill{} dim.\ $\leq n$. Alors, quitte à passer à une extension radicielle finie de $K$, on \uncertain{pourra} trouver un \add{\ill{}} morphisme surjectif \add{birationnel} et propre
\[
X' \longrightarrow X ,
\]
avec $X'$ lisse. \struck{\ill{}} La suite spectrale de \uncertain{descente} de Deligne donne
\[
H^*(X_{\bar K}, \mathbf{Q}_\ell) \Longleftarrow E_2^{pq} = H^p\bigl(\alpha \longmapsto H^q((X'/X)^{\alpha+1}_{\bar K}, \mathbf{Q}_\ell)\bigr) .
\]
\struck{Dans} \struck{Cette suite spectrale} \struck{Nous supposons \ill{} que l'\ill{} \ill{} $H^q$ pour \ill{} $q < i$, \ill{} \ill{} \ill{} sans restriction sur la dimension de $X$} L'hyp. de récurrence \ill{} \ill{} \ill{} \uncertain{signifie} \ill{} \struck{$E_2^{pq} \to E_\infty$} l'action de $I$ sur $E_2^{pq}$ est quasi-unipotente quand $q < i$. \struck{Pour} \struck{$q < i$, on doit prendre} Parmi ces termes $E_2^{pq}$ \ill{} $p + q = i$, le seul \ill{} \ill{} \ill{} \ill{} \struck{s'il est} si l'action de $I$ \ill{} quasi-unipotente \ill{} \ill{} $E_2^{0,i} \subset H^i(X'_{\bar K}, \mathbf{Q}_\ell)$. \struck{Mais} On \ill{} \ill{}

\note{Le bas de la page, de la proposition 9.6 à la fin, est barré de deux grandes croix, et chaque ligne y est en outre raturée ; seuls des fragments se lisent.}
\struck{\ill{} \ill{} \ill{} $K^{p^{-\infty}}$ des schémas de \ill{} dimension \ill{}. \textbf{Prop. 9.6.} \ill{} le th. de monodromie pour les \ill{} $X_K$ de t.f., \ill{} \ill{} pour $i \leq n$, \ill{} $H^i(X_{\bar K}, \mathbf{Q}_\ell)$ \ill{} \ill{} \ill{} \ill{} $X_K$ \ill{} \ill{} \ill{} \ill{} $i = n$, \ill{} $X_{\bar K}$ \ill{} lisse \uncertain{quasi-projectif}, et \ill{} \ill{} extension radicielle $K'$ de $K$. Alors \ill{} \ill{} \ill{} $H^n(X_{\bar K}, \mathbf{Q}_\ell)$ pour \ill{} $X$ \ill{} \ill{} \ill{} \ill{} la clôture \ill{} $K^{p^{-\infty}}$, admet une résolution des singularités.}

\page{49}
\note{p.~48 de l'auteur.}
\textbf{\struck{Corollaire} \add{Proposition} 9.6.} Supposons la résolution des singularités \add{(ordinaire)} sur $K^{p^{-\infty}}$ en \uncertain{toute} dimension, et que pour tout $X$ \struck{\ill{}} projectif et lisse sur $K$, \struck{\ill{}} l'action de $I$ sur $H^i_*(X_{\bar K}, \mathbf{Q}_\ell)$ soit quasi-unipotente, et \uncertain{kif-kif} après extension finie radicielle de $K$. Alors \struck{\ill{}} pour tout $X_K$ de t.f. sur $K$, l'action de $I$ sur les $H^i(X_{\bar K}, \mathbf{Q}_\ell)$ est quasi-unipotente.

Il suffit de \uncertain{conjuguer} \struck{\ill{}} \add{\ill{} quiproquo} 9.4 et 9.4~c), en procédant par récurrence sur $n$.
\note{« 9.4 » corrige un premier renvoi raturé ; le second porte « §~9.4 c) » ou « 9.4 c) ».}

\struck{\ill{}} 9.7. \emph{Cas des} $H^i(X_{\bar K}, \mathbf{Q}_\ell)$, pour $i \leq 1$.

Le cas de $i = 0$ est trivial et laissé au lecteur (en fait, $\pi$ y opère de façon \emph{quasi-triviale}).

Pour le cas $i = 1$, on \struck{\ill{}} peut, quitte à une extension finie de $K$, supposer que les composantes irréductibles réduites de $X_K$ sont géom. intègres. Supposons $X$ réduit, \ill{} \ill{} \ill{} \ill{} \ill{} \ill{}, $X'$ sa normalisée, $X'' = X' \times_X X'$, $X''' = X' \times_X X' \times_X X'$ \add{(\ill{} \ill{}, \uncertain{voir} SGA~4~VIII)} : \struck{\ill{} \ill{}} la suite spectrale de descente finie donne
\[
0 \to H^1\bigl(\alpha \mapsto H^0((X'/X)^{\alpha+1}_{\bar K}, \mathbf{Q}_\ell)\bigr) \to H^1(X_{\bar K}, \mathbf{Q}_\ell) \to \operatorname{Ker}\bigl(H^1(X'_{\bar K}, \mathbf{Q}_\ell) \overset{\partial}{\rightrightarrows} H^1(X''_{\bar K}, \mathbf{Q}_\ell)\bigr) .
\]
\note{Une première version de la ligne, raturée : $0 \to \operatorname{Ker}\bigl(H^0(X''_{\bar K}) \overset{\partial}{\rightrightarrows} H^0(X'''_{\bar K})\bigr)$. Sous le premier terme, une accolade : « \ill{} d'un sous-\ill{} de $H^0(X''_{\bar K}, \mathbf{Q}_\ell)$ ».}
Cela nous ramène à voir que l'action de $I$ sur \struck{\ill{}} $H^1(X'_{\bar K}, \mathbf{Q}_\ell)$. \emph{D}\supplied{onc} \ill{} \ill{} \uncertain{ramène} au

\page{50}
\note{p.~49 de l'auteur.}
cas de $X$ géom. normal. \struck{La suite \ill{} \ill{} $\mathcal{U} \ill{} \ill{}$} On peut \uncertain{supposer} de plus $X$ géom. intègre (\struck{\ill{} \ill{} \ill{}} \uncertain{les} composantes connexes l'étant). Si $U$ est un ouvert non vide de $X$, \ldots{} $H^1(X_{\bar K}, \mathbf{Q}_\ell) \hookrightarrow H^1(U_{\bar K}, \mathbf{Q}_\ell)$, donc on \struck{peut} \ill{}

\note{Le passage qui suit est encadré, barré de traits croisés, et chaque ligne y est en outre raturée.}
\struck{supposer $X$ affine, donc quasi-projectif. Soit $\hat X$ \ill{} \ill{} de $X$, \ill{} \ill{} \ill{} \ill{} donc quasi-projectif. Quitte à faire une extension \add{finie} radicielle sur $K$, \ill{} $X \subset \bar X$ \ill{} $\bar X$ projectif, \ill{} \ill{} que \ill{} \ill{} \uncertain{remplaçant} $X$ par $U$. \ill{} \ill{} \ill{} géom. normal. \uncertain{Notons} \ill{} \ill{} \ill{} \ill{}, $U$ un \ill{} \ill{} \ill{} \uncertain{dans} \ill{} \ill{}, \ill{} \ill{} $V$ de \ill{}. Alors \ill{}}

\struck{$H^1(\ill{}$} \struck{\ill{} \ill{} \ill{} \ill{} \ill{} \ill{} $\hat X$ \ill{} $U$, qui \ill{} \uncertain{supposé} affine.} \add{et lisse.} \struck{Soit \ill{} une clôture projective \add{$\hat X$} de $U$, qui \ill{}} \struck{\ill{} \ill{} $W_{K'/K}$ (\ill{} quitte à faire \ill{} une extension radicielle finie) géom. normale.} \struck{\ill{} $H^1(U, \mathbf{Z}/\ell^\nu\mathbf{Z})$ \ill{} \ill{} \ill{} $X$, \ill{} \ill{} \ill{} sur $X$, \ill{} \ill{} \ill{} \ill{} des \ill{} de $X - U$.}
\note{Plusieurs fragments de cette zone sont des ajouts de sa main au-dessus de lignes biffées, et ont été biffés à leur tour ; l'ordre de lecture n'est pas sûr.}

Une propriété \uncertain{d'extension} transcendante pour des $V$ \uncertain{« de \ill{} géométrique »} \ill{} \ill{} \ill{} \ill{} \ill{} \ill{} $X$ est de dim.\ $1$. \struck{Alors le calcul explicite} \add{\ill{} plus \ill{}} \struck{\uncertain{montre}} \ill{} le modèle projectif \uncertain{régulier} $\hat X$ de $X$ est lisse. (\ill{} La dualité nous ramène \ill{} \ill{} $H^1_!(X_{\bar K}, \mathbf{Q}_\ell)$, la suite exacte \ill{} \ill{} \ill{} \ill{} de $H^1_!(\hat X, \mathbf{Q}_\ell) = H^1(\hat X, \mathbf{Q}_\ell)$. \ill{} est donc ramené au cas d'une courbe \struck{projective} \add{\emph{lisse}} sur $K$. A \uncertain{expliciter} \ill{} \ill{})

\textbf{Proposition 9.8.} \ldots
\note{L'énoncé de la proposition 9.8 est laissé en blanc, marqué seulement de pointillés.}

\section*{Exposé II. Cas d'action \struck{\ill{}} \add{\uncertain{risquable}} du groupe fondamental}
\note{Titre de la p.~52. Au-dessous, entre parenthèses : « (Une variante du th. de monodromie) » ; un « II » est porté au-dessus de « fondamental ». La page précédente, p.~51, est blanche.}

\page{52}
Nous donnons \struck{dans le} \struck{les actions} présent exposé un résultat de nature analogue au théorème de monodromie de l'exposé précédent, \struck{\ill{}} à titre de nouvelle illustration de la méthode qui y était employée. Il ne servira plus dans la suite de ce Séminaire.

1. \emph{Énoncé du résultat}, \struck{\ill{}} \emph{corollaires}.

Écrivons le résultat \struck{essentiel} \add{principal du présent exposé}, \uncertain{c'est} le suivant :

\textbf{Théorème 1.1.} Soient $k$ un corps, \struck{$\bar k$ une clôture algébrique, $k$ \ill{} \ill{} \ill{} \ill{} \ill{} \ill{} tel que le corps des \ill{} déduit de $k$} \struck{\ill{} \ill{} \ill{}} (ou non absolu), $X$ un schéma sur \struck{$k$} \add{$S = \operatorname{Spec} k$}, \struck{\ill{}} de type fini (ou non absolu), normal, localement de type fini, géométriquement connexe, $\xi$ un point géométrique \add{de $X$}, d'où (SGA~1~VIII) une suite exacte
\[
(1.1.1) \qquad 1 \to \pi_1(\bar X, \xi) \longrightarrow \pi_1(X, \xi) \longrightarrow \pi_1(S, \xi) \longrightarrow 1 ,
\]
avec \struck{\ill{}} $\bar X = X \otimes_k \bar k$, $\bar k$ étant la clôture séparable de $k$ dans $k(\xi)$, (et $\xi$ désignant aussi les points géométriques de $S$ resp. $\bar X$ déduits canoniquement de $\xi$). \struck{(\ill{}) $I(\ell)$ le \ill{} \ill{}} \struck{Considérons le plus grand} Soient $\ell$ un nombre premier distinct de la caractéristique \struck{\ill{}} de $k$, \struck{\ill{}} $\pi$ un sous-groupe ouvert de $\pi_1(X,\xi)$, $\pi_0$ son image dans $\pi_1(S,\xi)$, $I = \pi \cap \pi_1(\bar X, \xi)$, de sorte qu'on a une suite exacte
\[
(1.1.2) \qquad 1 \longrightarrow I \longrightarrow \pi \longrightarrow \pi_0 \longrightarrow 1 .
\]
Soit $I(\ell)$ la composante $\ell$-primaire du plus grand quotient profini abélien $I_{\mathrm{ab}}$ de $I$, \struck{Alors} d'où \add{par passage au quotient une extension}
\[
(1.1.3) \qquad 1 \longrightarrow I(\ell) \longrightarrow \pi' \longrightarrow \pi_0 \longrightarrow 1 .
\]
Alors \struck{\ill{}} $I(\ell)$ est un module de type fini sur $\mathbf{Z}_\ell$, et $[\pi', I(\ell)]$ \struck{est \ill{}} est d'indice fini dans $I(\ell)$, en d'autres termes \struck{\ill{}} le plus grand quotient profini de $I(\ell)$ sur lequel $\pi_0$ opère trivialement est un groupe \emph{fini}.
\note{Les notations « $\bar X$ » et « $\bar k$ » rendent des lettres surlignées sur la page ; l'exposé lui-même est la suite du précédent, qui était l'exposé sur le théorème de monodromie.}

\page{53}
\note{p.~II~2 de l'auteur.}
Ce théorème \uncertain{sera} démontré dans les sections suivantes.

\note{Le passage qui suit est encadré et barré en croix.}
\struck{\ill{} \ill{} \ill{} \ill{} au cas où $X$ est un ouvert d'une courbe algébrique projective et lisse sur $k$, on peut \ill{} \struck{\ill{} \ill{} \ill{} \ill{}} l'énoncé résultant de l'hyp. de Riemann pour les courbes algébriques sur un corps fini, démontrée par A.~Weil \ill{} dans ce cas.}

Nous \ill{} dans le présent n° donnons quelques conséquences. Notons que la suite exacte (1.1) implique une suite exacte
\[
(I^{\mathrm{ab}})_{\pi_0} \longrightarrow \pi^{\mathrm{ab}} \longrightarrow \pi_0^{\mathrm{ab}} \longrightarrow 1
\]
\struck{\ill{}} et en passant aux composantes $\ell$-primaires
\[
I(\ell)_{\pi_0} \longrightarrow \pi^{\mathrm{ab}}(\ell) \longrightarrow \pi_0^{\mathrm{ab}}(\ell) \longrightarrow 1 ,
\]
d'où \ill{} concluant :

\textbf{Corollaire 1.2.} L'homomorphisme \add{surjectif} canonique
\[
\pi^{\mathrm{ab}}(\ell) \longrightarrow \pi_0^{\mathrm{ab}}(\ell)
\]
\struck{\ill{}} a un \emph{noyau fini}.

\struck{Utilisant I~6.1, \ill{} \ill{} \ill{} de 1.1 \ill{}}

\textbf{Corollaire 1.3.} Soit
\[
u : \pi_1(X,\xi) \longrightarrow \operatorname{Aut}_{\mathbf{Q}_\ell}(E)
\]
une représentation $\ell$-adique \struck{\ill{}}, telle que la représentation induite sur \add{$\pi_1(\bar X, \xi)$} \struck{\ill{}} soit \struck{\emph{abélienne}} \add{\emph{essentiellement résoluble}}, \struck{ou plus \ill{}} \struck{\emph{géométriquement} \ill{} \ill{} de Riemann \ill{}} (i.e. \uncertain{qu'un} sous-groupe ouvert de $I$ dans l'action soit \struck{abélien} \add{résoluble}). Alors la représentation induite sur $I$ est \emph{essentiellement unipotente}, \struck{\ill{}} (\uncertain{donc} essentiellement triviale si elle est semi-simple).
\note{Les notations $\pi$, $\pi_0$ et $I$ de cette page sont celles de 1.1 ; le renvoi « (1.1) » est sur la page ; « I~6.1 » renvoie à l'exposé précédent.}

\page{54}
\note{p.~II~3 de l'auteur.}
En effet, \struck{\ill{}} l'hypothèse \struck{\ill{} \ill{} \ill{}} \add{\ill{} \ill{}} \uncertain{implique} qu'il existe un sous-groupe \struck{\ill{} \ill{} \ill{} \ill{}} ouvert $\pi$ de $\pi_1(X,\xi)$ tel que la représentation induite sur $I = \pi \cap \pi_1(\bar X, \xi)$ soit \ill{} \ill{} petit, \ill{} \ill{} \uncertain{supposer} abélienne. Prenons $U$ \ill{} (\ill{} $I'$) qui \struck{\ill{} \ill{} \ill{} \ill{}} \ill{} l'image de \struck{$I$} dans $\operatorname{Aut}_{\mathbf{Q}_\ell}(E)$ \ill{} un pro-$\ell$-groupe. Donc la représentation \add{de $\pi$} provient d'une représentation de $\pi'$, \struck{\ill{}} l'extension (1.3) de $\pi_0$ par $I(\ell)$. On conclut par I~6.1.

\textbf{Corollaire 1.4.} Avec les notations de 1.3, supposons \struck{\ill{}} l'action de $\pi_1(X,\xi)$ \uncertain{commutative} \struck{\ill{}} \add{\emph{essentiellement}} abélienne. Alors la représentation induite sur $I$ est essentiellement triviale, \struck{\ill{}} \struck{\ill{} \ill{} \ill{} \ill{} surjectif \ill{}} \add{i.e.} il existe un sous-groupe \struck{\ill{}} ouvert \struck{\ill{} $X' \to X$,} $U$ de $\pi$ tel que la restriction de la représentation \add{\ill{}} associée à $\pi$ \struck{\ill{}} se factorise par \ill{} \struck{\ill{}}
\[
\pi_0 = \operatorname{Im}\bigl(\pi \to \pi_1(S,\xi)\bigr) .
\]
\struck{Cet résultat \ill{} \ill{}} On prend \struck{\ill{}} d'abord $\pi$ tel que la restriction à $\pi$ soit abélienne, \struck{\ill{} \ill{}} et que \struck{\ill{} \ill{}} l'image de $\pi$ dans $\operatorname{Aut}_{\mathbf{Q}_\ell}(E)$ soit un pro-$\ell$-groupe. Alors la représentation de $\pi$ \ill{} représentation provient d'une représentation de $\pi^{\mathrm{ab}}(\ell)$, et \struck{il suffit d'\ill{} que} \add{\ill{}} 1.2 \ill{} \ill{} qu'il existe un sous-groupe ouvert $U$ de $\pi^{\mathrm{ab}}(\ell)$ tel que $U \to \pi_0^{\mathrm{ab}}(\ell)$ soit injectif ; donc la représentation restreinte à $U$ provient d'une représentation de son image dans $\pi_0^{\mathrm{ab}}(\ell)$. Il suffit

\page{55}
\note{p.~II.4 de l'auteur.}
donc de \struck{considérer} \add{prendre} l'image inverse $V$ de $U$ dans $\pi$, la restriction de \uncertain{la} $u$ à $V$ provient d'une représentation de l'image de $V$ dans $\pi_0^{\mathrm{ab}}(\ell)$, donc d'une repr. de l'image de $V$ dans $\pi_0$, cqfd.

\note{Dans la marge gauche, ce diagramme ; les flèches obliques sont tracées d'un trait plus fin, et l'on a gardé leur sens tel qu'il se lit.}

\begin{tikzcd}[column sep=small, row sep=small]
 & V \arrow[r, hook] \arrow[d] \arrow[dl] & \pi \arrow[d] \\
\operatorname{Im} V \arrow[r, hook] \arrow[dr] & U \arrow[r, hook] \arrow[d, "\wr"] \arrow[dr] & \pi^{\mathrm{ab}}(\ell) \arrow[d] \\
 & U \arrow[r, hook] & \pi_0^{\mathrm{ab}}(\ell)
\end{tikzcd}

\note{La flèche de $\operatorname{Im} V$ vers le $\pi_0$ voisin s'écrit sur la page « $\operatorname{Im} V \hookrightarrow \pi_0$ » ; on l'a rendue par la seule flèche oblique que le dessin porte vers le $U$ du bas. Un trait oblique joint aussi $\pi$ à $U$.}

On va donner une application de ces résultats à la théorie des V.A.

\textbf{Théorème 1.5.} Soient $k$, \struck{$A$}, \add{\struck{$X$, $\xi$} comme dans 1.1,} \add{$k$ un corps sép. clos, $X$ schéma loc. de type fini, mais sép. clos} \struck{\ill{} \ill{} \ill{} \ill{} normal et \ill{}}, $\xi$ un point géométrique de $X$, \struck{\ill{}} $\ell$ un nombre premier distinct de la caractéristique de $k$, $A$ un schéma abélien sur $X$, \add{\emph{géométriquement intègre}}. Supposons que la représentation de $\pi_1(\bar X, \xi)$ dans $T_\ell(A_\xi)$ soit \struck{abélienne} \add{essentiellement} résoluble. Alors cette représentation est essentiellement \struck{unipotente} \add{triviale}, \struck{donc essentiellement triviale si elle est semi-simple (*)}. \struck{Donc \ill{} \ill{}} et il existe un revêtement étale fini surjectif $X' \to X$, \struck{\ill{} extension finie séparable $k'$ de $k$, \ill{}} \struck{$X' \to \operatorname{Spec}(k')$ \ill{}} \add{\ill{}} un schéma abélien $B$ sur $k$, \struck{tel que $A_{X'}$ \ill{} \ill{} radicielle} et une \emph{isogénie} \add{\emph{radicielle}} $B_{X'} \to A_{X'}$. \struck{qui soit \ill{}} En particulier, si $k$ est de car. nulle, \struck{\ill{}} $A_{X'}$ est isogène à $B_{X'}$ (i.e. est un schéma abélien constant).

Prouvons d'abord l'assertion que la représentation de $\pi_1(X,\xi)$ dans $T_\ell(A_\xi)$ est \add{ess. unipotente} \struck{\ill{} \ill{}}. \struck{Donc} Soit $U$ un voisinage \ill{} \ill{} \ill{} de $\xi$,

\note{Au pied de la page, sous un trait, l'appel (*) : « (*) Il y a lieu de conjecturer que (sans hypothèse \ill{}) l'action de $\pi_1(X,\xi)$ sur $T_\ell(A_\xi)$ est toujours semi-simple », le tout biffé. La colonne $A$ / $X$ / $S$ en marge gauche figure la tour $A \to X \to S$.}

\page{56}
\note{p.~II.5 de l'auteur.}
alors, \struck{$X$} $U$ étant normal, $\pi_1(U,\xi) \to \pi_1(X,\xi)$ est surjectif, donc il \struck{est \ill{} \ill{} \ill{}} suffit de prouver l'assertion pour $U,\xi$ et $A|U$. Cela nous permet de supposer $X$ affine, donc de t.f. sur $k$.

\note{En marge gauche, ce carré :}

\begin{tikzcd}[column sep=small, row sep=small]
A_0 \arrow[r, no head] \arrow[d, no head] & A \arrow[d, no head] \\
X_0 \arrow[r, no head] \arrow[d, no head] & X \arrow[d, no head] \\
S_0 \arrow[r, no head] & S
\end{tikzcd}

On sait alors que $X$, $A$ proviennent de $S_0 = \operatorname{Spec} k_0$, $k_0$ un sous-corps de t.f. de $k$, $X_0$ schéma de t.f. sur $S_0$, $A_0$ schéma abélien sur $X_0$. Alors, \struck{$\xi$ \ill{} \ill{} \ill{}} \struck{$X$} $X \to X_0$ étant fid. plat et $X$ normal, $X_0$ est normal. De plus, $\xi$ définit un pt géométrique de $X_0$, d'où par conséquent \struck{\ill{} \ill{}} la suite exacte du type (1.1.1)
\[
1 \to \pi_1(\bar X_0, \xi) \longrightarrow \pi_1(X_0, \xi) \longrightarrow \pi_1(S_0, \xi) \to 1 .
\]
\struck{De plus, \ill{} Soit si \ill{} \ill{}} \ill{} \ill{} $\bar k_0$ est la clôture séparable de $k_0$ dans $k$, on a $\bar X_0 \simeq X_0 \otimes_{k_0} \bar k_0$, donc $X_0 \simeq \bar X_0 \otimes_{\bar k_0} k$. On a un hom.
\[
(1.5.1) \qquad \pi_1(X,\xi) \longrightarrow \pi_1(\bar X_0, \xi)
\]
qui est un \uncertain{homomorphisme} surjectif, \struck{\ill{}} (\uncertain{vérif}). Par ces conditions, la représentation naturelle de $\pi_1(X,\xi)$ dans $T_\ell(A_{0\xi})$ \add{$= T_\ell(\bar A_\xi)$}, et la représentation induite via $\pi_1(X,\xi) \to \pi_1(\bar X_0, \xi) \hookrightarrow \pi_1(X_0, \xi)$ \ill{} \ill{} autre que celle qu'il s'agit d'étudier \add{$A$}, \struck{\ill{}} \add{ess. résoluble} (resp. \struck{\ill{}} \ill{} \ill{}) la surjectivité de (1.5.1), dire qu'elle \add{est} \uncertain{ess. unipotente} signifie que elle induite sur $\pi_1(\bar X_0, \xi)$ l'est. \struck{Donc} On conclut donc par 1.3.
\note{Dans la formule ${X_0 \simeq \bar X_0 \otimes_{\bar k_0} k}$, le premier membre est bien écrit $X_0$ sur la page, alors que l'argument semble demander $X$.}

\page{57}
\note{p.~II.6 de l'auteur.}
Il existe donc un ss-groupe ouvert $\pi' \hookrightarrow \pi_1(X,\xi)$ sur lequel la repr. est unipotente. Mais $\pi'$ définit un rev. étale connexe $X'$ de $X$, et $\xi$-pointé, tel \ldots

\note{Le passage qui suit, jusqu'à « est triviale », est barré de traits obliques ; la phrase ci-dessus, ajoutée au-dessus de lui, semble le remplacer. Dans la marge gauche, une variante (cf. infra).}
\struck{Nous savons maintenant que la représentation envisagée de $\pi_1(X,\xi)$ \uncertain{étant} aussi semi-simple, donc, étant unipotente, elle est triviale sur un sous-groupe \add{ouvert $U$} \struck{\ill{}} de $\pi_1(X,\xi)$. Ce sous-groupe définit un revêtement \add{connexe $\xi$-pointé} $X'$ de $X$, et par construction la représentation de $\pi_1(X',\xi')$ dans $T_\ell(A_{\xi'})$ est triviale.}
\marginal{\ill{} que la représentation naturelle de $\pi_1(X',\xi') = \pi'$ dans $T_\ell(A_{X'}(\xi'))$ soit unipotente. On \ill{} la démonstration, \ill{} \ill{} \uncertain{déjà} que $\pi_1(X,\xi)$ opère unipotentement.}

Il suffit alors d'appliquer \uncertain{l'adaptation suivante d'}un résultat bien connu de Lang--\uncertain{Néron}, \struck{\ill{}} utilisant Mordell--Weil (cf. I~\ill{} \ill{}~4.4) : \struck{\ill{} \ill{} \ill{}}

\textbf{Lemme 1.5.1.} Si l'action de $\pi_1(X,\xi)$ sur $T_\ell(A_\xi)$ est \struck{\ill{}} unipotente, \add{\struck{\ill{} \ill{}}} \struck{si} cette action est triviale, et $\exists$ schéma abélien $B$ sur $k$, et une isogénie radicielle $B_X \to A$.

\textbf{Corollaire 1.6.} \struck{Sous les hyp.} Si pour \struck{un} $\ell \neq \mathrm{car}\, k$, l'action de $\pi_1(X,\xi)$ sur $T_\ell(A_\xi)$ est \add{ess.\ résoluble} \struck{\ill{} \ill{}} \struck{\ill{} \ill{}} ($\Leftrightarrow$ \struck{\ill{}} essentiellement triviale) il en est de même pour tout autre $\ell' \neq \mathrm{car}\, k$.

\textbf{Corollaire 1.7.} Supposons \add{car $k = 0$,} \struck{$A$} l'action de $\pi_1(X,\xi)$ sur $T_\ell(A_\xi)$ soit unipotente et semi-simple, \struck{\ill{}} \add{\uncertain{entier}} $n \geq 3$, \struck{\ill{} \ill{} \ill{}} \add{$\ell \neq$ car $k$,} et supposons qu'il existe un \struck{\ill{} \ill{}} \add{\uncertain{tel}} tel que la représentation \struck{de} induite \add{de} $\pi_1(X,\xi)$ sur ${}_nA(\xi)$ soit triviale, i.e. que ${}_nA$ soit un revêtement \uncertain{induit} groupes \emph{triviaux} de $X$. Alors \struck{\ill{} 1.5 \ill{} \ill{}} \struck{\ill{} \ill{} \ill{}, si car $k = 0$,} $A$ est un schéma abélien \emph{constant} sur $X$.

En effet, \struck{il} \add{la} \uncertain{théorie} de la descente \struck{\ill{} $\pi_1(X,\xi)$ opère \ill{} $T_\ell(A_\xi)$ \ill{}} \struck{l'intersection d'un groupe \emph{fini}. Or un} \add{\uncertain{géométrique} et} la structure discrète de $\operatorname{Aut}_k(B)$ impliquent que
\note{La marge porte aussi : « $B_X \to A$ isogénie radicielle », relié par une flèche au lemme 1.5.1.}

\page{58}
\note{p.~II.7 de l'auteur.}
$A$ est décrit par une représentation
\[
\Gamma \longrightarrow \operatorname{Aut}(B)
\]
où $\Gamma$ est le groupe de Galois de $X'/X$. L'hypothèse implique que $\Gamma$ opère trivialement sur ${}_nB(k)$, et on sait que (comme $n \geq 3$) cela implique que l'action de $\Gamma$ sur $B$ est triviale, d'où la conclusion.

\textbf{Remarque 1.8.} \struck{Le théorème 1.5 \ill{}} En \uncertain{changeant}, on voit que 1.7 \uncertain{reste} vrai \ill{} \ill{} tel que l'on suppose $k$ \struck{\ill{}} \add{\uncertain{séparablement}} clos.

\struck{L'}\emph{\uncertain{énoncé}} 1.5 est de nature purement géométrique (bien que sa démonstration soit de nature arithmétique). Voici une variante arithmétique de 1.5 :

\textbf{Théorème 1.9.} Soit $K$ un corps de type fini \struck{\ill{}}, $\bar K$ une clôture séparable de $K$, $\pi = \operatorname{Gal}(\bar K / K)$, $A_K$ un schéma abélien sur $K$. Supposons l'action de $\pi$ sur $T_\ell(A(\bar K))$ \struck{abélienne} \add{essentiellement résoluble}. Alors il existe une extension finie \add{$k' \subset \bar K$} \struck{\ill{}} du corps \ill{} \ill{} premier $\mathfrak{k}$, \struck{et un} un schéma abélien $B$ sur $k'$, et une isogénie radicielle \struck{$B_{K'} \simeq A$}
\[
B \otimes_{k'} K' \simeq A \otimes_K K' ,
\]
$K' = K k'$. Lorsque car $K = 0$, \struck{\ill{} \ill{} \ill{}} \add{\ill{}} et que de plus $\ell$ \ldots
\note{La fin de l'énoncé est récrite par-dessus ; on lit : « l'action de $\pi$ sur $T_\ell(A(\bar K))$ est ess.\ abélienne (*), ou tout au moins \ill{} \ill{}, \ill{} $k'$ soit \struck{\ill{} \ill{}} un corps de nombres, \ill{} \ill{} supposer que $B$ \struck{\ill{} \ill{}} provienne d'un schéma abélien sur $R$ \struck{\ill{}} (l'anneau des entiers de $k'$). »}

\note{Au pied, sous un trait : « (*) Hypothèse \uncertain{toujours} \uncertain{vérifiée} \ill{} \struck{\ill{}} \ill{} \ill{} (\ill{} \ill{}) : \ill{} \uncertain{conjecture} que l'action est \uncertain{très} ½simple ! » — le « ½ » est sa notation pour « semi- ».}

\page{59}
\note{p.~II.8 de l'auteur.}
1.10. \emph{Démonstration.} \struck{En effet, \ill{} \ill{} \ill{} 1.5.} \add{\uncertain{Existence} de $B$.} Soit $k$ la \struck{clôture} \add{\uncertain{fermeture}} algébrique du corps premier dans $K$, et soit $X$ un schéma \add{intègre} \struck{de type fini} affine \struck{\ill{} $K$ \ill{} une extension régulière} \add{de corps des fonctions $K$}, \struck{\ill{}} sur $k$, tel que $A_K$ provienne d'un schéma abélien sur $X$. \struck{Alors $X$ est \ill{}} On \uncertain{peut} supposer $X$ \struck{géom.} normal (car $K$ régulier), \add{\ill{}} et quitte à remplacer $k$ par une extension finie, il est loisible de supposer $X$ géom. intègre (N.B. $k$ est parfait, étant une extension finie du corps premier). Alors l'extension $\bar K$ de $K$ définit un pt géométrique $\xi$ de $X$, et l'hyp. sur $A_K$ signifie aussi que la représentation de $\pi_1(X,\xi)$ dans $T_\ell(A_\xi)$ est ess. \struck{abélienne} \add{résoluble}. En vertu de 1.3, \struck{\ill{}} il \ldots

\begin{tikzcd}[column sep=small, row sep=small]
X \arrow[r, no head] \arrow[d, no head] & X' \arrow[d, no head] \\
k \arrow[r, no head] & k'
\end{tikzcd}

\note{Le passage qui suit, jusqu'à « une isogénie radicielle », est encadré et barré en croix ; une version de remplacement est écrite par-dessus et autour.}
\struck{\ldots{} il existe un revêtement étale \add{surjectif} $X'$ de $X$ \ill{} \ill{} une extension finie $k'$ de $k$, \ill{} \ill{} que $T_\ell(A_\xi)$ \ill{} \ill{} \ill{} d'une représentation de $\pi_1(\operatorname{Spec}(k'),\xi)$, i.e.\ \ill{} \ill{} faisceau $T_\ell(A_{X'})$ \ill{} un faisceau \uncertain{localement} \ill{} \ill{} sur $X'$. En vertu du résultat déjà cité de Lang--Néron, il s'ensuit que \ill{} \ill{} \ill{} \ill{} $B$ sur $k'$ \ill{} \ill{} et une isogénie radicielle}
\[
B_{X'} \longrightarrow A_{X'} ,
\]
\add{\ill{} $X'$ de $X$ et \ill{} \uncertain{extension} finie $k'$ \ill{} \ill{} \uncertain{radicielle} : \ill{} $B_{k'} \otimes_{k'} X \to A_{X'}$ isogénie, \ldots{} \ill{} \ill{} la situation.}
\note{Le statut de ce dernier ajout par rapport au passage barré n'est pas sûr ; il semble pris dans le cadre.}

En car.\ quelconque. En car.\ $p > 0$, comme $\operatorname{Gal}(\bar k'/k') \simeq \hat{\mathbf{Z}}$ est abélien, on conclut que l'action de $\operatorname{Gal}(\bar K/K)$ sur $T_\ell(A(\bar K))$ est \struck{\ill{} \ill{} \ill{}} \ill{} fait \emph{abélienne}.

\struck{Remarque 1.10. \ill{} \ill{} \ill{} \ill{} \ill{} \ill{}}
\marginal{Cas de l'action supposée ess.\ abélienne aux pts de car.\ $\neq \ell$.}

1.11. \emph{Borne inférieure} \struck{\ill{}} \ldots{} Si $\bar k'$ est la clôture sép. de $k'$ dans $\bar K$, il est \ill{} \ill{} \ldots{} que l'on \uncertain{sait} \uncertain{bien} \uncertain{connu}.
\note{Le titre « 1.11 » et la ligne qui le suit sont reliés par un trait à la note marginale ; la ligne en pointillés est de lui.}

\page{60}
\note{p.~II.9 de l'auteur.}

\begin{tikzcd}[column sep=small, row sep=small]
\operatorname{Gal}(\bar K'/K') \arrow[r] \arrow[d] & \operatorname{Gal}(\bar k'/k') \\
\operatorname{Gal}(\bar K/K) &
\end{tikzcd}

est \uncertain{donc} \uncertain{injectif} puisque \ill{} \ill{} deux sous-groupes ouverts. Donc \struck{l'hyp.} \add{\ill{}} dire que \ill{} $A_K$, ou $A_{K'}$, ou $B$ sur $k'$, donnent lieu à des actions ess.\ \struck{abéliennes} \struck{\uncertain{résolubles}} \add{abéliennes}, revient au même (Ceci est vrai chaque fois que \ill{} une extension de type fini\ldots{}) ! On en conclut p.\ ex.\ déjà une réciproque \struck{\ill{}} de l'énoncé \add{dans le cas d'un corps fini}, lorsque car $K > 0$ : s'il existe $B/k'$ avec $B \otimes_{k'} K'$ isogène à $A \otimes_K K'$, alors l'action de \struck{$\pi_1(K/K')$ sur} $\operatorname{Gal}(\bar K/K)$ \struck{sur} sur $T_\ell(A(\bar K))$ est \struck{\ill{}} ess.\ abélienne : en effet, celle de \struck{\ill{}} $\operatorname{Gal}(\bar k'/k') \simeq \hat{\mathbf{Z}}$ sur $T_\ell(B(\bar k))$ l'est évidemment !
\note{À la dernière ligne, un $\hat{\mathbf{Z}}$ biffé précède $\operatorname{Gal}(\bar k'/k')$, et l'égalité « $\simeq \hat{\mathbf{Z}}$ » est un ajout au-dessus de la ligne.}

Dans le cas car.\ $K = 0$ nous ne savons montrer que \ill{} \ill{} \ill{} 1.9, \ill{} \ill{} \ill{} \ill{} où \ill{} $K$ est déjà un corps de nombres. Dans ce cas, \struck{d'après \ill{}} le résultat est bien connu. Par \struck{\ill{}} \struck{\ill{} \ill{}} \ill{} \ill{} \ill{} \struck{\ill{}} \ill{} et \ill{} \ill{} critère de Néron--\add{Ogg}--Šafarevič \add{(\ill{} \ill{} \ill{}, cf.\ aussi Exp.~\ill{})} on est ramené à voir que si $R$ est l'anneau des entiers de $K$, \add{\ill{} \ill{} idéal maximal \ill{} de $R$} \struck{l'action du} \add{\ill{} du} groupe d'inertie $I \subset \pi = \operatorname{Gal}(\bar K/K)$ \ill{}, \ill{} relativement : \ill{} est quasi-\struck{unipotente} \add{triviale} (de sorte qu'après une extension finie convenable $K'$ de $K$, \ldots

\marginal{\uncertain{décidé} \ill{} \ill{}}
\note{Le bas de la page, à partir de « Dans le cas car.\ $K = 0$ », est marqué d'un long trait vertical dans la marge. La phrase est interrompue en bas de page et se poursuit au-delà de ce lot. Le crochet « [ ] » et la référence « ([\ill{} \ill{}~2]) » qui précèdent « cf. aussi Exp. » sont de sa main et restent illisibles. Le nom « Šafarevič » est rendu tel qu'il se lit, « Chafarevitch » n'étant pas exclu.}

\end{document}
